\chapter{Introducción}
\label{sec:introduccion}

En el centro de la Vía Láctea reside una fuente de radio compacta y brillante, Sagitario~A*, cuya naturaleza se ha asociado con un agujero negro supermasivo. Durante las últimas décadas, el seguimiento astrométrico de las estrellas del cúmulo~S, que orbitan esta fuente a distancias de apenas milésimas de pársec, ha permitido usar su dinámica como una sonda gravitacional directa del campo central \cite{gillessen2017monitoring}. El monitoreo sostenido durante más de dos décadas ha reconstruido órbitas keplerianas cerradas en torno a un cuerpo oscuro de alrededor de cuatro millones de masas solares, confinado en una región tan compacta que difícilmente admite otra interpretación que la de un agujero negro supermasivo. En este régimen de campo fuerte, el movimiento de dichas estrellas se convierte en un banco de pruebas natural para las predicciones de la relatividad general \cite{parsa2017investigating}.

La protagonista de este programa observacional ha sido la estrella S2, con un período orbital cercano a dieciséis años y un pericentro muy próximo a la fuente. En el año 2020, la GRAVITY Collaboration detectó en su órbita la precesión de Schwarzschild \cite{gravity2020schwarzschild}, es decir, el avance del pericentro que la relatividad general predice como primera corrección al problema kepleriano. Esta medición constituye la primera evidencia del comportamiento relativista de la órbita, aunque todavía dentro de un tratamiento que ignora la rotación de la fuente. Poco después, en el año 2022, la Event Horizon Telescope Collaboration (en adelante EHTC) publicó la primera imagen de la sombra de Sagitario~A* \cite{eht2022sgraI}, resultado que se considera prueba directa de la existencia de un agujero negro supermasivo en el centro de nuestra galaxia y que, además, abre la puerta a la caracterización de sus parámetros físicos.

El movimiento de una estrella ligada admite una jerarquía de aproximaciones sucesivas. En el límite newtoniano la órbita es una elipse cerrada que no precesa. La relatividad general introduce, como primera corrección, el avance del pericentro —la precesión de Schwarzschild ya mencionada—, que depende únicamente de la masa de la fuente y no de su estado de rotación \cite{einstein1916fundamentos}. El orden siguiente incorpora los efectos del momento angular del agujero negro: la rotación arrastra consigo a los sistemas inerciales vecinos y rompe la simetría entre las órbitas prógradas y las retrógradas, añadiendo a la trayectoria una precesión del propio plano orbital \cite{bardeen1972rotating}. Es en este orden donde el parámetro de rotación deja su huella sobre la dinámica estelar, y por ello la detección de la precesión de Schwarzschild marca el punto de partida natural del paso siguiente.

Confirmada la hipótesis del agujero negro supermasivo, el parámetro que resta por determinar es el \emph{parámetro de rotación} $a$, restringido al intervalo $a\in(-1,1)$ en unidades geométricas, donde $a>0$ corresponde a órbitas prógradas y $a<0$ a órbitas retrógradas. Su magnitud no es un mero detalle: en la \emph{métrica de Kerr} \cite{kerr1963gravitational}, que describe el campo de un agujero negro en rotación, el valor de $a$ determina el tamaño de la ergosfera —la región donde ningún observador puede permanecer en reposo— y desplaza el radio de la órbita circular estable más interna (ISCO) según el sentido del movimiento \cite{bardeen1972rotating}. La rotación, por tanto, reconfigura la estructura misma de las órbitas admisibles alrededor de la fuente.

Junto con la imagen, la EHTC argumentó a favor de un valor prógrado, y desde entonces la literatura se ha decantado por un parámetro de rotación elevado. A partir de los datos del EHTC y de observaciones simultáneas en rayos~X, se ha estimado un valor de $0.9\pm0.6$ mediante el método de \emph{outflow} \cite{daly2023new}; se ha establecido una cota inferior en torno a $0.8$ a partir de la radiación del plasma \cite{yfantis2023fitting}; y la comparación con modelos magnetohidrodinámicos favorece un valor cercano a $0.94$ \cite{eht2024sgraVIII}.

Todas estas determinaciones, sin embargo, parten de la emisión del plasma en las cercanías del horizonte, y no de la dinámica orbital de las estrellas. Descansan, por lo tanto, sobre modelos magnetohidrodinámicos del disco de acreción y sobre supuestos acerca de la geometría de emisión, cuyas incertidumbres se propagan al valor inferido. La dinámica orbital ofrece, frente a ello, un enfoque complementario e independiente: sondea la métrica de manera directa, sin la mediación del complejo comportamiento del plasma, caracterizando el parámetro de rotación a partir de la evolución de las órbitas estelares en la métrica de Kerr. La determinación puntual del parámetro por esta vía, no obstante, dista de ser trivial. Por ello, en lugar de perseguir una estimación puntual, nos proponemos \emph{acotarlo superiormente}, aprovechando la estabilidad de las órbitas de las estrellas del cúmulo~S como criterio físico: dado que la rotación reconfigura la estructura orbital, las trayectorias ligadas y regulares que se observan a lo largo de décadas restringen los valores del parámetro compatibles con lo medido.

Para avanzar hacia este objetivo recorremos varias etapas. Primero construimos un marco teórico, con revisión bibliográfica del estado del arte, sobre la determinación del parámetro de rotación de Sagitario~A*. Sobre esa base desarrollamos y publicamos RelatiPy, una librería de Python de código abierto orientada a la investigación en relatividad general con énfasis en geodésicas, que integra geodésicas de tipo tiempo en las métricas de Schwarzschild y de Kerr mediante integradores cuya conservación de las cantidades de movimiento —la energía $E$, el momento angular $L_z$ y la constante de Carter $Q$— hemos verificado de forma sistemática. La librería incorpora además la solución analítica de las geodésicas ligadas en la métrica de Kerr \cite{fujita2009analytical} y una caché persistente de trayectorias que facilita la reutilización de los resultados por parte de cualquier investigador. Con esta herramienta implementamos métodos numéricos para la determinación de las posiciones y velocidades de las estrellas y el ajuste de las observaciones, y estudiamos la evolución de órbitas elípticas, con un rango de parámetros similar al de las órbitas del cúmulo~S, tanto computacional como analíticamente.

El desarrollo de estas ideas exige, ante todo, describir el campo gravitacional de un agujero negro en rotación y las trayectorias que las estrellas siguen en él. Comenzamos, pues, por la métrica de Kerr.
